
\section*{Среднее арифметическое и геометрическое трёх чисел}
\begin{theorem}[AM–GM для трёх чисел]
Для \(a,b,c\ge0\) выполнено \(a+b+c\ge3\sqrt[3]{abc}\), причём равенство достигается тогда и только тогда, когда \(a=b=c\).
\end{theorem}
\begin{proof}Положим \(x=\sqrt[3]a,y=\sqrt[3]b,z=\sqrt[3]c\). Тождество
\[
x^3+y^3+z^3-3xyz=\tfrac12(x+y+z)\bigl((x-y)^2+(y-z)^2+(z-x)^2\bigr)
\]
даёт неотрицательность левой части. Если \(x+y+z=0\), все три числа нулевые. Иначе равенство возможно ровно при \(x=y=z\). Это равносильно \(a=b=c\).\end{proof}
